% WHERE A PARAMETER CHARACTER IS DOUBLED.
%
% A token list SHOWN anywhere writes a parameter character twice, so that
% what is written could be read back to give the same list again. Measured
% against pdftex in four places at once:
%
%   \everypar={A#B}, an error inside it   <everypar> A##
%   \def\m{C#D}, an error inside it       \m ->C##
%   \toks0={G#H} \showthe\toks0           > G##H.
%   \show of a macro, and \meaning        already doubled here
%
% This engine wrote one # in the first three. It is the same rule the
% fourth already followed: a list is written the way \show writes one
% WHEREVER it is written.
%
% What is NOT doubled is a macro's PARAMETER text: `\def\p#1x#2{...}' shows
% as `\p #1x#2->...'. Those are not parameter characters in a list but the
% marks of where an argument goes, and each is written once with its number
% after it.
%
% trip line 121 puts a mark holding om#ega back to be read, and every error
% under it shows the text.
\catcode`\{=1 \catcode`\}=2 \catcode`\#=6
\tracingonline=1 \hsize=100pt
\message{[A a parameter character in a token parameter]}
\everypar={A#B}
\indent \par
\message{[B one in a macro's body]}
\def\m{C#D}
\m \par
\message{[C one in a token register shown]}
\toks0={G#H}
\showthe\toks0
\message{[D but a parameter text is written as it stands]}
\def\p#1x#2{A#1B##C#2}
\show\p
\message{[done]}
\end
